\documentclass[a4paper,11pt]{ltjsarticle}
\usepackage[haranoaji]{luatexja-preset}
\usepackage{amsmath,amssymb}
\usepackage[top=12mm,bottom=10mm,left=14mm,right=14mm]{geometry}
\usepackage{array}
\usepackage{tikz}
\usetikzlibrary{calc}
\pagestyle{empty}
\setlength{\parindent}{0pt}
\newlength{\qcol}\setlength{\qcol}{0.47\textwidth}
\newlength{\qgap}
\newlength{\anscolw}\setlength{\anscolw}{0.30\textwidth}
\newlength{\ansh}
\newcommand{\Q}[2]{\parbox[t]{\qcol}{\raggedright (#1)\hspace{0.6em}$\displaystyle #2$}}
\newcommand{\QR}[4]{\Q{#1}{#2}\hfill\Q{#3}{#4}\par\vspace{\qgap}}
\newcommand{\anscell}[1]{\rule[-\ansdrop]{0pt}{\ansh}(#1)}
\newlength{\ansdrop}

\begin{document}
{\large\bfseries 計算問題　30}\hfill 名前(\underline{\hspace{32mm}})\par\vspace{3mm}
■（1）〜（16）の計算をしなさい。（17）〜（21）は方程式を解きなさい。\par\vspace{3mm}
\setlength{\qgap}{5.4mm}
\QR{1}{\dfrac{8}{3}\div5+\dfrac{8}{3}}{2}{8(3x-2)-(5x-8)}
\QR{3}{(-12)\times(-5)}{4}{8\div3\times(-15)}
\QR{5}{\dfrac{x-8}{3}\times15}{6}{16-(-3)\div\left(-\dfrac{3}{5}\right)}
\QR{7}{(24a-15)\div3}{8}{(8a+3)+(5a-8)}
\QR{9}{\dfrac{1}{8}(24x-40)+\dfrac{1}{8}(40x+24)}{10}{(-24)\div(+3)}
\QR{11}{(3x-8)-(8x+5)}{12}{9\times\{3-(-5+1)\times2\}}
\QR{13}{-0.8-(-5.3)-1.8}{14}{-8^2+(-3)^2-5}
\QR{15}{(-15)+(-6)-(-10)}{16}{\dfrac{8}{3}(3x-24)-\dfrac{5}{8}(8x+64)}
\QR{17}{8x-3=29}{18}{8(x+3)=9x+18}
\QR{19}{0.8x+0.3=4.3}{20}{\dfrac{3}{4}(x+2)=\dfrac{1}{2}(x-3)}
\vspace{1mm}
\Q{21}{\dfrac{8x-3}{5}-\dfrac{x+8}{5}=5}\par\vspace{3mm}
\setlength{\ansh}{9.5mm}\setlength{\ansdrop}{6mm}%
\begin{center}\begin{tabular}{|p{\anscolw}|p{\anscolw}|p{\anscolw}|}\hline
\anscell{1} & \anscell{2} & \anscell{3} \\\hline
\anscell{4} & \anscell{5} & \anscell{6} \\\hline
\anscell{7} & \anscell{8} & \anscell{9} \\\hline
\anscell{10} & \anscell{11} & \anscell{12} \\\hline
\anscell{13} & \anscell{14} & \anscell{15} \\\hline
\anscell{16} & \anscell{17} & \anscell{18} \\\hline
\anscell{19} & \anscell{20} & \anscell{21} \\\hline
\end{tabular}\end{center}

\end{document}
